% ============================ CONCLUSION ==============================
\chapter*{Conclusion générale et perspectives}
\addcontentsline{toc}{chapter}{Conclusion générale et perspectives}
\markboth{CONCLUSION GÉNÉRALE}{}

Ce projet de fin d'études a permis de concevoir et de réaliser, au sein de la
\textbf{BTK -- Banque Tuniso-Koweïtienne}, une application web de gestion des
agences bancaires adossée à un volet décisionnel. La solution centralise la
gestion des agences, des employés, des clients et des objectifs, formalise les
opérations sensibles au moyen de circuits de validation, et restitue aux
responsables des indicateurs d'aide à la décision via un tableau de bord embarqué
et des rapports Power BI.

Sur le plan technique, ce travail a été l'occasion de mettre en pratique les
compétences acquises durant la formation : modélisation UML, développement
Jakarta EE (Servlets, JSP, EJB, JPA), conception de bases de données Oracle et
mise en œuvre d'une chaîne décisionnelle (vues analytiques et reporting Power BI),
au croisement du développement logiciel et de la Business Intelligence.

Plusieurs perspectives d'évolution peuvent être envisagées :
\begin{itemize}[noitemsep]
  \item le chiffrement (hachage) des mots de passe et le renforcement de la
        sécurité ;
  \item l'automatisation de l'alimentation d'un entrepôt de données (ETL) pour
        historiser les indicateurs ;
  \item l'enrichissement par des analyses prédictives (prévision des objectifs) ;
  \item l'exposition d'une API REST pour l'interopérabilité avec d'autres
        systèmes ;
  \item le développement d'une application mobile de consultation des indicateurs.
\end{itemize}
